% =====================================================================
% Paper 4 -- first full draft, targeted at Phys. Rev. E
%
% STATUS NOTES FOR THE AUTHOR (delete before submission)
%  * Every number in this draft is taken from hyperuniformity_1024.json
%    (1024^2, c0 = 0.30, 4 seeds, 2000 steps) and hyperuniformity_512.json.
%    Nothing is rounded from memory.
%  * References checked against the published records on 17 Aug 2026. Note:
%    the Yeung and Furukawa TITLES supplied in the final edit round did not
%    match the published articles (the volumes and pages did); the titles below
%    are the published ones. Tomita and Ezoe page ranges confirmed by the
%    author. Torquato et al. (2026), De Luca et al. (2024) and Tomita (1991)
%    are retained because Secs. I, III B and V B cite them; removing them
%    would leave undefined citations and would drop the closest prior art from
%    the reference list.
%    Yeung PRL 61, 1135 (1988) proves the BOUND delta >= 4, not delta = 4;
%    the text now says so, because our own alpha_c = 3.81 +/- 0.08 sits
%    2.4 sigma BELOW that bound and a referee will notice.
%  * Sec. IV D and Sec. V C now report the mass-autocorrelation measurement.
%    Note that this CHANGED a claim: an earlier draft of Sec. IV C called the
%    separation of the residual "future work". It is no longer future work,
%    and the sentence was removed rather than left standing.
%  * NOVELTY WARNING, READ BEFORE SUBMITTING. Torquato, Kim, Klatt, Car and
%    Steinhardt, Phys. Rev. X 16, 011042 (2026), generalise hyperuniformity to
%    weighted particle systems and show explicitly that a hyperuniform point
%    pattern need not be hyperuniform when weighted, and vice versa. That is
%    the general form of the conceptual claim this paper makes. Secs. I and
%    V A have been rewritten to cite it and to position our contribution as
%    the conserved-dynamics instance plus the variance budget, but the author
%    must read that paper in full and re-examine whether more of the framing
%    needs to change. Similarly De Luca et al. (2024) and Tomita (1991)
%    establish the k^4 phase-ordering result more directly than Yeung and
%    Furukawa do, and Wilken et al., PRX 13, 031014 (2023), may already have
%    measured droplet point patterns in a phase-separating system: read it.
%  * The mineralogical setting is deliberately reduced to one paragraph in
%    Sec. II A, per the decision to target PRE. Papers 1-3 carry it.
% =====================================================================

\documentclass[aps,pre,twocolumn,superscriptaddress,floatfix,longbibliography]{revtex4-2}

\usepackage{graphicx}
\usepackage{amsmath}
\usepackage{amssymb}
\usepackage{url}

\begin{document}

\title{Conserved dynamics make the solute hyperuniform, not the domains}

\author{Ruqing Chen}
\email{ruqing@hotmail.com}
\affiliation{GUT Geoservice Inc., Montr\'eal, Qu\'ebec, Canada}

\date{\today}

\begin{abstract}
Hyperuniformity---the anomalous suppression of long-wavelength density
fluctuations---is increasingly used to classify spatial patterns in physical,
biological, and geological systems. A pattern is normally classified by
locating its constituent objects and measuring the structure factor of the
resulting point set. We show that for a system with a locally conserved order
parameter this procedure can give the wrong answer, because the conserved
quantity and the counted objects are not the same thing. Using Cahn--Hilliard
coarsening of a two-phase alloy free energy on periodic grids up to
$1024^2$, we measure the spatial order of the solute-rich domains against
matched complete-spatial-random, hard-core, and perturbed-lattice nulls. The
domains are locally more regular than a hard-core gas (Clark--Evans
$R = 1.73$, $z = +73$, against $1.42$ for random sequential adsorption). The
concentration field is hyperuniform, with $S_c(k) \sim k^{3.81 \pm 0.08}$. The
point pattern of domain centroids is not: $S(k \to 0)$ settles on a plateau
near $0.14$ and the number variance grows with window area, and this survives
a doubling of the box while the bias of the estimators themselves disappears.
Weighting the same points by domain solute mass restores hyperuniformity,
$S(k \to 0) = 0.021$. In common windows the solute-mass variance is a factor
of eight below what independent domains of the observed size distribution
would produce, and the domain count and the mean domain size anticorrelate
ever more strongly with window size. Number fluctuations are absorbed by
compensating size fluctuations: permuting the domain masses over the same
positions returns the budget exactly to its independent-domain baseline,
locating the entire effect in the coupling between mass and position. Nearest-neighbour regularity and long-range
order must be measured separately, and so must counts and conserved
quantities.
\end{abstract}

\maketitle

% =====================================================================
\section{Introduction}
% =====================================================================

A point pattern is \emph{hyperuniform} if its large-scale density fluctuations
are anomalously suppressed: the structure factor vanishes in the
long-wavelength limit, $S(k) \to 0$ as $k \to 0$, equivalently the variance of
the number of points in a window of radius $R$ grows more slowly than the
window's volume---in two dimensions like its perimeter, $\sigma^2(R) \sim R$,
rather than like its area, $\sigma^2(R) \sim R^2$
\cite{TorquatoStillinger2003,Torquato2018}. The concept has proved unusually
portable. Disordered hyperuniform patterns have been identified in avian
photoreceptor mosaics \cite{Jiao2014}, in jammed packings, in active and
driven matter, and the diagnostic is now routinely applied to patterns in
biology and the geosciences.

Portability has a cost. Hyperuniformity is defined for a density, but it is
almost always \emph{measured} on a point pattern obtained by locating discrete
objects---cells, particles, grains, depressions---and treating their centroids
as the density's carriers. The step from the density to the point pattern is
rarely examined. It is exact when every object carries the same weight, and it
is a genuine approximation otherwise.

Two recent results frame the question we ask here. On the measurement side,
Chen \cite{Chen2026} applied the full protocol---a local diagnostic, the
structure factor, the number variance, and matched nulls---to fields of
submarine pockmarks, seabed depressions left by escaping fluid. The
nearest-neighbour spacing of one field was more regular than a hard-core gas,
the kind of local order that invites a hyperuniform interpretation; yet every
field was long-range super-Poissonian, with $S(k \to 0)$ between $2.5$ and
$19.4$, and the hyperuniform region of the order plane was empty. The
methodological conclusion was that a local statistic cannot substitute for a
long-range one. That system, however, is one in which long-range order simply
does not exist, so the point-pattern step was never put under stress.

On the theory side, systems governed by a locally conserved order parameter
are the natural place where long-range order \emph{does} exist. The
Cahn--Hilliard equation \cite{CahnHilliard1958} is a continuity equation: a
composition fluctuation of wavelength $2\pi/k$ can only be built by
transporting solute across that distance, so long-wavelength modes are
kinematically frozen while short-wavelength ones grow. The resulting
suppression of the small-$k$ structure factor of the concentration field is
long established: Yeung \cite{Yeung1988} proved that a scaling coarsening
pattern with conserved dynamics must have $S_c(k) \sim k^{\delta}$ with
$\delta \ge 4$, numerical work found $\delta \simeq 4$
\cite{Furukawa1989}, Tomita \cite{Tomita1991} gave the scaling argument that
explains it, and De Luca \emph{et al.} \cite{DeLuca2024} recently extended
that argument to active and non-constant-mobility cases, again obtaining
$4$. We claim no novelty whatever for this, and reproduce it here only as an
internal check on the simulation.

The question we ask is what happens to the \emph{point pattern} in such a
system. Mass conservation constrains the solute; it does not constrain how
many domains hold it. A region can meet its solute budget with a few large
domains or with many small ones, and nothing in the dynamics forbids the
trade. Whether the centroids inherit the field's hyperuniformity is therefore
an empirical question, and the answer bears directly on every study that
classifies a conserved system by counting objects.

That the answer can be negative is, in general, now known. Torquato
\emph{et al.} \cite{TorquatoWeighted2026} have generalised hyperuniformity to
particles carrying weights---charges, dipoles, Voronoi volumes---and shown
that a hyperuniform arrangement need not remain hyperuniform when weighted,
and that a non-hyperuniform one may become so; related weighted point
constructions have been examined in an ecological setting
\cite{Ezoe2025}. The general statement that positions and weights are
different objects is therefore established, and we do not claim it.

What this paper adds is the conserved-dynamics instance and its mechanism.
First, we exhibit a system in which the weight is not an arbitrary internal
degree of freedom but the conserved quantity itself, and show that the field,
the centroid count pattern, and the mass-weighted centroid pattern of one
simulation fall into different classes---so that the discrepancy is not a
curiosity of exotic weights but the generic situation whenever discrete
objects share out a conserved density. Second, we quantify the difference: a
variance budget, built from an identity and three measured inputs with no free
parameter, shows the observed solute rigidity to be a factor of eight beyond
what independent domains of the observed size distribution would give, with
the shortfall supplied by an anticorrelation between domain count and mean
domain size that deepens monotonically with scale. Third, we draw the
methodological consequence for the many studies that classify such systems by
locating and counting objects.

% =====================================================================
\section{Model and methods}
% =====================================================================

\subsection{The system}

We integrate the Cahn--Hilliard equation
\begin{equation}
  \frac{\partial c}{\partial t} = \nabla \cdot \left[ M \nabla
  \left( \frac{\partial f}{\partial c} - \kappa \nabla^2 c \right) \right],
\end{equation}
for a conserved composition $c(\mathbf{r}, t)$ on a periodic $n \times n$
grid, with a semi-implicit pseudospectral scheme. The homogeneous free energy
$f(c)$ is a double well with a solute-poor and a solute-rich minimum,
calibrated in earlier work for a solid solution of a minor element in a
sulphide host; the specific chemistry is immaterial here and is reported
elsewhere. What matters for the present argument is only that $f$ has two
minima, that the initial composition lies inside the spinodal, and that the
dynamics conserve $\int c \, d\mathbf{r}$ locally rather than merely globally.
All runs use $c_0 = 0.30$, off-critical, so that the minority phase forms
discrete domains rather than a bicontinuous structure and centroids are
well defined.

The gradient coefficient $\kappa$ is fixed by requiring the fastest-growing
spinodal wavelength to be $12$ grid cells, and the time step by requiring the
fastest-growing mode to be amplified by $2\%$ per step. The absolute length
scale is therefore set by numerical resolvability and not by the material;
accordingly, every quantity reported below is dimensionless---ratios to a
Poisson reference, scaling exponents, and correlation coefficients---and none
depends on that choice.

The principal ensemble is four independent runs at $n = 1024$, coarsened
isothermally for $2000$ steps, at which point each configuration contains
$2736 \pm 20$ domains with a mean spacing of $19.6$ cells. A second ensemble
of six runs at $n = 512$, carried to $4000$ steps, is used for the
finite-size comparison of Sec.~\ref{sec:finitesize}. We stop the large
ensemble at the stage that maximises the accessible wavevector range rather
than at the latest attainable time: coarsening reduces the domain count, and
with it the statistical quality of precisely the long-wavelength band the
paper is about.

\subsection{From field to point pattern}

Domains are the connected components of $c > c_{\rm thr}$, with $c_{\rm thr}$
the midpoint between robust percentiles of the two phase compositions.
Components are labelled on the torus, so that a domain straddling the periodic
boundary is one domain and not two. Each domain contributes one point, the
centroid of its excess composition $c - c_{\rm thr}$, computed as a circular
mean so that the wrap is handled correctly, and one weight, its solute mass
$m_i = \sum (c - c_{\rm thr})$ over the component. Components smaller than
eight cells are discarded as unresolved ripples. Moving $c_{\rm thr}$ across
the central $40\%$ of the composition gap changes the exponents reported below
by less than $0.05$.

\subsection{Diagnostics}

\emph{Local order.} The Clark--Evans index \cite{ClarkEvans1954}
$R = \bar d_{\rm NN} / [1/(2\sqrt{\rho})]$, with $R < 1$ clustered, $R = 1$
random, $R > 1$ regular. On a torus the minimum-image distance is exact and no
edge correction arises---an advantage over field surveys, where boundary bias
is the dominant systematic.

\emph{Long-range order.} The structure factor is evaluated on the exact
reciprocal lattice of the periodic box,
\begin{equation}
  S(\mathbf{k}) = \frac{1}{N} \left| \sum_{j=1}^{N}
  e^{-i \mathbf{k} \cdot \mathbf{r}_j} \right|^2 ,
  \qquad \mathbf{k} = 2\pi (n_x, n_y)/L ,
\end{equation}
with no window function and hence none of the low-$k$ leakage that
complicates bounded surveys, and radially averaged in logarithmic shells. For
an uncorrelated set of $N$ points on this lattice $\langle S \rangle = 1$ for
every $\mathbf{k} \neq 0$. The weighted version,
$S_w(\mathbf{k}) = |\sum_j w_j e^{-i\mathbf{k}\cdot\mathbf{r}_j}|^2 /
\sum_j w_j^2$, is normalised so that uncorrelated positions again give unity
whatever the weights; with $w_j = m_j$ it is a spectrum of the solute rather
than of the domains. The number variance $\sigma^2(R)$ is taken over disks
placed at random on the torus. We fit the low-$k$ exponent $\alpha$ in
$S \sim k^{\alpha}$ between the smallest available wavevector and
$0.45 k_{\rm peak}$, and quote it with its standard error.

\subsection{Two constraints that must not be read as results}
\label{sec:constraints}

First, the scheme conserves the mean composition exactly and the box is
periodic, so $S_c(k = 0) = 0$ identically. That is arithmetic, not physics;
only $k > 0$ is fitted or plotted. Second, a configuration with $N$ points in
a periodic box has a number variance driven to zero as the window approaches
the whole box, for the same trivial reason. Windows are therefore restricted
to $R \le L/6$ (Sec.~\ref{sec:constraints}), and---importantly---the complete-spatial-random null is
generated at \emph{fixed} $N$, so that it carries the identical constraint and
the comparison remains fair.

\subsection{Nulls and validation}

Each configuration is compared against complete spatial randomness (CSR) at
matched $N$; a random-sequential-adsorption (RSA) hard-core process with
exclusion matched to the observed closest approach, representing steric
exclusion without organisation; and a Gaussian-perturbed square lattice as a
known hyperuniform reference. Passed through the identical pipeline, CSR
returns $R = 0.99$, $S(k \to 0) = 1.14$, and a number-variance slope of
$2.02$, against the ideal $1$, $1$, and $2$; the perturbed lattice returns
$\alpha = +2.02$ and slope $0.93$. The estimators are therefore doing what
they are supposed to, and the null results reported below are not artefacts of
the pipeline.

% =====================================================================
\section{Results}
% =====================================================================

\begin{table*}[t]
\caption{The three objects and the nulls, at $1024^2$ after $2000$ steps
($N = 2736$ domains per configuration, four realisations). $R$ is the
Clark--Evans index, $H = S(k\to0)/S(k_{\rm peak})$, and the last column the
log--log slope of the number variance fitted above $1.4$ domain spacings.}
\label{tab:main}
\begin{ruledtabular}
\begin{tabular}{lccccc}
 & $R$ & $S(k\to0)$ & $\alpha$ & $H$ & $\sigma^2$ slope \\
\hline
Domain centroids (counts) & 1.73 & 0.143 & $-0.14(4)$ & 0.089 & 1.80 \\
Same points, mass-weighted & --- & 0.021 & $+0.40(9)$ & 0.014 & --- \\
Concentration field & --- & --- & $+3.81(8)$ & --- & --- \\
\hline
CSR (fixed $N$) & 0.99 & 1.14 & $-0.05$ & --- & 2.02 \\
RSA hard core & 1.42 & 0.24 & $+0.24$ & 0.20 & 1.79 \\
Perturbed lattice & 1.45 & 0.017 & $+2.02$ & 0.008 & 0.93 \\
\end{tabular}
\end{ruledtabular}
\end{table*}

\subsection{The domains are locally more regular than a hard-core gas}

Table~\ref{tab:main} collects the results discussed in this section.

The centroid pattern has $R = 1.73$ ($z = +73$), well above both CSR ($0.99$)
and the RSA hard-core null ($1.42$) at the same density [Fig.~\ref{fig:one}]. Local order
is thus real and is not merely exclusion: the domains are spaced more evenly
than hard discs of the same closest approach would be, as expected when each
domain draws solute from, and so depletes, its neighbourhood.

\begin{figure}[t]
\includegraphics{fig_1}
\caption{(a) Solute-rich domains (grey scale) and their mass-weighted
centroids (circles), a $384^2$ crop of one $1024^2$ configuration after
$2000$ steps. (b) Order plane: Clark--Evans index against the
long-wavelength structure factor. The domains are locally regular, exceeding
the hard-core null, but their count spectrum sits well above the hyperuniform
limit; weighting the same points by solute mass moves them down by nearly an
order of magnitude. Grey squares are the three Mediterranean pockmark fields
of Ref.~\cite{Chen2026}, shown for scale comparison only.}
\label{fig:one}
\end{figure}

\subsection{The field is hyperuniform}

The concentration field has $S_c(k) \sim k^{\alpha_c}$ with
$\alpha_c = 3.81 \pm 0.08$ over more than a decade in $k$ [Fig.~\ref{fig:two}(a)],
recovering the established behaviour of conserved coarsening
\cite{Yeung1988,Furukawa1989,Tomita1991,DeLuca2024}. We report it as a check
on the simulation, not as a finding.

The fitted value lies nominally below the bound $\delta \ge 4$ of
Ref.~\cite{Yeung1988}, by about $2.4$ standard errors of the fit. We do not
read this as a violation. The bound holds asymptotically for a scaling state,
whereas the fit is taken over a finite window that ends at $0.35 k_{\rm peak}$,
where curvature towards the peak is already depressing the local slope;
restricting the fit to the lowest half-decade returns a steeper exponent with
a correspondingly larger error. The quoted uncertainty is that of the
regression alone and does not include this systematic.

\begin{figure}[t]
\includegraphics{fig_2}
\caption{(a) Structure factor of the concentration field, with a $k^4$ guide.
(b) Spectra of the point patterns and nulls. The two centroid curves are the
\emph{same points}, differing only in whether each point carries unit weight
or its domain's solute mass. Every series is distinguished by line style and
marker as well as colour, so the figure is legible in monochrome. The
dot-dashed vertical line marks the coarsening peak $k_{\rm peak}$.}
\label{fig:two}
\end{figure}

\subsection{The point pattern is not}

The centroid pattern is not hyperuniform. Its structure factor falls well
below the Poisson line but then flattens, settling on a plateau
$S(k \to 0) = 0.143$ with a low-$k$ exponent $\alpha = -0.14 \pm 0.04$
[Fig.~\ref{fig:two}(b)], and its number variance grows as $\sigma^2 \sim R^{1.80}$, close
to the area-like exponent $2$ and far from the perimeter-like $1$ returned by
the perturbed lattice [Fig.~\ref{fig:three}(a)]. By the criterion that defines the class,
the pattern is not hyperuniform; by the ratio of the plateau to the peak,
$H = 0.089$, it is not effectively hyperuniform either. It is, however,
strongly suppressed: its number variance lies a factor of four to five below
CSR at every window size, and below the hard-core null as well.

\subsection{Doubling the box does not move the plateau}
\label{sec:finitesize}

A plateau at the smallest accessible wavevector is exactly what a box too
small to resolve the true limit would produce, so we repeated the measurement
at $n = 512$ and $n = 1024$, halving $k_{\rm min}$ from $0.0123$ to $0.0061$
cell$^{-1}$. At the same coarsening stage the two spectra coincide over their
shared range, and the additional half-decade of low $k$ opened by the larger
box shows the same plateau [Fig.~\ref{fig:three}(b)]. The exponents agree within their
errors, $-0.20 \pm 0.08$ at $512^2$ against $-0.14 \pm 0.04$ at $1024^2$.

The strongest form of this argument is asymmetric. Between the two box sizes
the \emph{estimators} improve markedly---the CSR control moves from
$(S, {\rm slope}) = (0.92, 1.85)$ to $(1.14, 2.02)$, converging on its
theoretical values as the finite-box bias is removed---while the data do not
move at all. A residual bias that vanishes for the null but leaves the
measurement unchanged cannot be what produces the measurement.

\begin{figure}[t]
\includegraphics{fig_3}
\caption{(a) Number variance against window radius, with slope-1 and slope-2
guides. (b) Finite-size check: the count spectrum in boxes of $512^2$ and
$1024^2$ at the same coarsening stage; dotted verticals mark the respective
$2\pi/L$. Doubling the box halves the accessible $k_{\rm min}$ and leaves the
plateau where it was, while over the same change the CSR control moves from
$(S, {\rm slope}) = (0.92, 1.85)$ to $(1.14, 2.02)$. The bias that vanishes
for the null is not what produces the measurement.}
\label{fig:three}
\end{figure}

\subsection{Weighting by solute mass restores hyperuniformity}

Weighting the identical set of points by domain solute mass drops the
long-wavelength plateau from $0.143$ to $0.021$ and turns the exponent
positive, $\alpha = +0.40 \pm 0.09$, with $H = 0.014$ [Fig.~\ref{fig:two}(b)]. Nothing
about the positions has changed. What has changed is which quantity is being
counted, and the two answers differ by nearly an order of magnitude.

\subsection{A negative result: filtering by size does not isolate the cause}
\label{sec:negative}

The obvious way to test whether polydispersity is responsible is to restrict
the analysis to domains of nearly equal mass. It does not work, and the way it
fails is instructive. Selecting a subset removes points, and random removal by
itself drives $S(k)$ towards the Poisson value---for independent thinning with
retention $p$, $S_{\rm thinned}(k) = 1 + p[S_{\rm full}(k) - 1]$---so the
comparison must be against a random subset of the same size, not against the
full pattern. Made that way, the test returns the wrong sign: retaining the
$30\%$ of domains closest to the median mass gives $S(k \to 0) = 0.963$
against $0.616$ for matched random thinning, and the discrepancy grows as the
selection is tightened (Table~\ref{tab:filter}).

The diagnosis is that the two subsets have \emph{identical} Clark--Evans
indices ($1.19$ versus $1.19$ at $q = 0.30$; $1.35$ versus $1.35$ at
$q = 0.50$). The excess is therefore not local: it is long-wavelength. Domain
mass is spatially autocorrelated on large scales, so a subset selected by mass
is also a subset selected by position, and the selection injects exactly the
kind of long-wavelength density variation the measurement is designed to
detect. A hard filter cannot separate polydispersity from the spatial
autocorrelation of the filtering variable. We report the result because it
motivates the analysis of the next section, which deletes nothing.

% =====================================================================
\begin{table}[t]
\caption{The negative filtering test of Sec.~\ref{sec:negative}. Retaining the
fraction $q$ of domains whose mass is closest to the median is compared with a
random subset of the same size drawn from the same configuration. The
filtered subsets have higher long-wavelength fluctuations, and identical local
order.}
\label{tab:filter}
\small
\begin{ruledtabular}
\begin{tabular}{c@{\hspace{4pt}}c@{\hspace{4pt}}c@{\hspace{4pt}}c@{\hspace{4pt}}c@{\hspace{4pt}}c}
$q$ & CV & $S(k\to0)$ & $S(k\to0)$ & $R$ & $R$ \\
 & kept & filtered & thinned & filtered & thinned \\
\hline
0.30 & 0.12 & 0.963 & 0.616 & 1.19 & 1.19 \\
0.50 & 0.21 & 0.599 & 0.587 & 1.35 & 1.35 \\
0.70 & 0.31 & 0.482 & 0.417 & 1.51 & 1.51 \\
\end{tabular}
\end{ruledtabular}
\end{table}

\section{The variance budget}
% =====================================================================

\subsection{An identity with no free parameter}

For a sum of $N$ masses drawn independently of one another and of $N$, the
variance of a compound random sum is
\begin{equation}
  \frac{{\rm Var}[M]}{\langle m \rangle^2}
  = \langle N \rangle \, {\rm CV}^2 + {\rm Var}[N] ,
  \label{eq:budget}
\end{equation}
with ${\rm CV}$ the coefficient of variation of the domain masses. Every
quantity on the right is measured: $\langle N \rangle$ and ${\rm Var}[N]$ in
the same windows in which $M$ is accumulated, and ${\rm CV} = 0.528$ from the
mass distribution. Equation~(\ref{eq:budget}) is therefore a prediction with no
adjustable parameter, and the residual
\begin{equation}
  C \equiv \frac{{\rm Var}[M]}{\langle m \rangle^2}
  - \langle N \rangle \, {\rm CV}^2 - {\rm Var}[N]
\end{equation}
is the part of the observed solute rigidity that independence cannot account
for.

The value of this construction is that ${\rm CV}$ is an input that the window
statistics do not already contain. It converts the qualitative statement
``size compensates count'' into a number: how much fluctuation independent
domains of the observed size distribution \emph{would} have produced.

\subsection{Independence is excluded by a factor of eight}

Table~\ref{tab:budget} gives the budget across window sizes and Fig.~4
displays it. At the widest window, $R = 170.7$ cells or about eight domain
spacings, independent domains would give a variance of $100.3$ in units of one
domain, of which $66.6$ comes from the size term and $33.8$ from the count
term. The measured value is $12.2$. The residual is $C = -88.1$: the system
cancels $88\%$ of the fluctuation that polydispersity alone would impose.

The cancellation is not a coincidence of one window size. The ratio of
measured to baseline falls monotonically from $0.38$ at $R \approx 1.4$
spacings to $0.12$ at $R \approx 8$ spacings [Fig.~\ref{fig:four}(b)]: the wider the
window, the more completely the compensation operates. Over the same range the
correlation between the domain count and the mean domain mass within the
window deepens monotonically from $-0.62$ to $-0.89$, while the spread of the
window-mean domain mass contracts from $0.221$ to $0.031$. Where there are
more domains, the domains are smaller.

\begin{table*}[t]
\caption{The variance budget, Eq.~(\ref{eq:budget}), in units of one domain.
${\rm CV}^2 = 0.279$ throughout. Windows narrower than $1.4$ domain spacings
are omitted; there the decomposition is dominated by discreteness. The last
column is the correlation between the domain count and the mean domain mass in
the same window.}
\label{tab:budget}
\begin{ruledtabular}
\begin{tabular}{ccccccc}
$R$ & $\langle N\rangle$ & ${\rm Var}[N]$ & $\langle N\rangle{\rm CV}^2$ &
baseline & ${\rm Var}[M]/\langle m\rangle^2$ & ${\rm corr}(N,\bar m)$ \\
\hline
 27.7 &   6.3 &  1.36 &  1.75 &   3.11 &  1.17 & $-0.62$ \\
 35.9 &  10.6 &  2.14 &  2.95 &   5.09 &  1.57 & $-0.68$ \\
 46.5 &  17.7 &  3.38 &  4.95 &   8.33 &  2.19 & $-0.74$ \\
 60.4 &  29.9 &  5.50 &  8.33 &  13.83 &  3.03 & $-0.79$ \\
 78.3 &  50.2 &  8.49 & 14.00 &  22.49 &  4.14 & $-0.83$ \\
101.5 &  84.4 & 13.75 & 23.56 &  37.31 &  5.95 & $-0.86$ \\
131.6 & 141.8 & 22.81 & 39.57 &  62.38 &  8.27 & $-0.88$ \\
170.7 & 238.6 & 33.76 & 66.57 & 100.33 & 12.24 & $-0.89$ \\
\end{tabular}
\end{ruledtabular}
\end{table*}
\begin{figure*}[t]
\includegraphics{fig_4}
\caption{The variance budget, Eq.~(\ref{eq:budget}), in units of one domain.
(a) What independent domains would produce, split into its size and count
terms, against the measured solute-mass variance. The single input not already
contained in the window statistics is the coefficient of variation of the
domain masses, ${\rm CV} = 0.528$; it is what makes the baseline a prediction
with no free parameter. (b) The surviving fraction of that baseline, and the
correlation between the domain count and the mean domain mass in the same
window. \emph{These two curves are not independent evidence}: where the solute
content of a window is nearly rigid, $\bar m \propto 1/N$ is close to an
identity and ${\rm corr}(N,\bar m) \to -1$ follows algebraically. The
correlation is shown because it locates the scale on which the compensation
sets in, not as a second measurement. Shading marks windows narrower than
$1.4$ domain spacings, where the decomposition is dominated by discreteness.}
\label{fig:four}
\end{figure*}

\subsection{What the budget does and does not establish}

Three limitations bound this analysis, and we state them explicitly because
each could otherwise be over-read.

First, the correlation ${\rm corr}(N, \bar m) \to -0.89$ and the smallness of
${\rm Var}[M]$ are \emph{not} independent pieces of evidence. If the solute
content of a window is nearly rigid, then $\bar m \propto 1/N$ is close to an
identity and a correlation approaching $-1$ follows algebraically. The
correlation is reported because it expresses the same fact in the language of
the mechanism and supplies the scale on which the compensation sets in, not
because it constitutes a second measurement.

Second, the residual $C$ mixes two effects: correlation between a domain's
mass and the local domain count, and correlation between the masses of
neighbouring domains. The decomposition itself does not separate them; the
measurements of Sec.~\ref{sec:attribution} do.

Third, the budget should not be inverted to predict the plateau height of
Sec.~III C. Doing so would require the additional assumption that domain
positions are determined by the mass field, which we have not established. The
real-space cancellation and the reciprocal-space plateau are two views of one
physical situation and are reported side by side, not derived from one
another.

\subsection{Attribution by mark permutation and mass autocorrelation}
\label{sec:attribution}

Two further measurements close the budget. Both treat the domain masses as
marks attached to the observed positions, and both use a mark permutation as
their null: the masses are randomly reassigned among the same domains, which
destroys any relation between mass and position while leaving the positions,
the mass distribution, and hence ${\rm CV}$ exactly as observed.

\emph{The baseline is not an approximation.} Recomputing the budget with
permuted marks returns residuals consistent with zero at every window size,
and a ratio of measured to baseline variance of $0.99$--$1.09$ against
$0.88$--$0.12$ for the real marks. Equation~(\ref{eq:budget}) therefore
describes these configurations exactly once the marks are made independent of
position, which is what it assumes; the entire residual $C = -88$ is
attributable to mass--position coupling and none of it to a failure of the
identity.

\emph{The coupling is with crowding, not with other masses.} The mark
covariance $c_{mm}(r) = \langle \delta m_i \delta m_j\rangle_r/\sigma_m^2$
departs from its permutation envelope only out to about two domain spacings
and lies inside it, indistinguishable from zero, from $\approx 50$ cells to
the largest separation measured, $232$ cells [Fig.~\ref{fig:five}(a)]. Domain
mass carries no large-scale spatial structure. The correlation between a
domain's mass and the number of neighbours within $r$, by contrast, reaches
$-0.68$ at one domain spacing and decays slowly to $-0.08$ at $128$ cells
[Fig.~\ref{fig:five}(b)], against a permuted value of $+0.01$ throughout.
Crowded domains are smaller. The compensation that makes the solute rigid is
therefore local and steric in origin rather than a large-scale organisation of
domain sizes.

% =====================================================================
\begin{figure*}[t]
\includegraphics{fig_mass_autocorrelation}
\caption{Domain mass treated as a mark on the point pattern, against a
mark-permutation null in which the masses are reassigned at random among the
same domains. (a) Mark covariance. The signal is confined to separations below
about two domain spacings; beyond $\approx 50$ cells it is inside the null
envelope out to $232$ cells, so domain mass has no large-scale spatial
structure. Values exceeding unity in the first two bins reflect the small
number of pairs available there, where the null envelope is correspondingly
wide, and should not be interpreted. (b) Correlation between a domain's mass
and the number of neighbours within $r$: crowded domains are smaller. This is
the coupling that supplies the residual of the variance budget.}
\label{fig:five}
\end{figure*}

\section{Discussion}
% =====================================================================

\subsection{Relation to weighted hyperuniformity}

The formal apparatus for what we measure here exists. Reference
\cite{TorquatoWeighted2026} defines weighted spectral densities and local
variances for particles carrying scalar or tensorial weights, and establishes
that the weighted and unweighted classifications are logically independent.
Our contribution is not that apparatus but a physical instance of it in which
the weight is fixed by a conservation law rather than chosen: the solute mass
of a coarsening domain. That case has, we believe, a status the general
framework does not confer on it. When the weight is a conserved density, the
weighted spectrum is not one of many possible characterisations---it is the
one the dynamics constrain, and the unweighted spectrum is the one experiment
usually reports. The gap between them is then not a mathematical possibility
but a systematic error with a known sign.

\subsection{A rule for classifying conserved systems}

The results support a short procedural rule for anyone applying
hyperuniformity diagnostics to a pattern of discrete objects. (i) Measure
local and long-range order separately; a high nearest-neighbour index is
compatible with area-like number variance, as it is here and as it was in the
pockmark fields \cite{Chen2026}. (ii) If the system has a conserved quantity,
measure both the count spectrum and the spectrum weighted by that quantity;
here they differ by a factor of seven and fall on opposite sides of the
classification boundary. (iii) Run the null through the identical pipeline and
report what it returns, because the credibility of a negative result rests
entirely on the estimator having been shown to work.

The rule applies well beyond the present model. Precipitate populations,
droplet emulsions, cluster distributions in reaction-diffusion systems, and
ecological patch mosaics are all conserved or quasi-conserved systems that are
routinely characterised by locating objects and counting them.

\subsection{The empirical case}

Domain point patterns in a phase-separating system have been measured
experimentally. In phase-separated DNA droplets, Wilken \emph{et al.}
\cite{Wilken2023} report structure factors and number variances whose
long-wavelength fluctuations do not vanish, and attribute the failure to
incomplete equilibration and to kinetic arrest of the coalescence process.
% This characterisation of Ref. [Wilken2023] follows the corresponding author's
% reading of that paper; check the wording against the original before
% submission.
The present results suggest that a further contribution is available without
invoking either: even in a fully passive, unarrested, isothermally coarsening
system, the count spectrum does not become hyperuniform, because the count is
not the conserved variable. Where droplet volumes are measurable alongside
positions, the weighted spectrum is the discriminating measurement, and our
prediction is definite---it should lie well below the unweighted one.

\subsection{Two scales, one lesson}

The pockmark analysis of Ref.~\cite{Chen2026} and the present work sit at
opposite ends of a very wide range of scales---seabed depressions tens of
metres apart, and nanometre-scale precipitates---and reach the same
methodological conclusion from opposite directions. There, local regularity
exceeding a hard-core null coexisted with strongly super-Poissonian long-range
behaviour, and the warning was that local order must not be read as long-range
order. Here, long-range order genuinely exists, and yet the standard
point-pattern measurement misses it, because the conserved quantity is not the
counted object.

The two systems also differ in a way worth stating. In the pockmark fields the
excess long-wavelength fluctuation was traced to substrate heterogeneity: the
inhomogeneous-Poisson nulls reproduced the observed $S(k \to 0)$, so the
long-range signal was a property of the geological background rather than of
the pattern-forming process. In the simulations there is no substrate, no
heterogeneity, and no external field, and the residual count fluctuation is
generated by the pattern-forming dynamics themselves. The lesson common to
both is that the long-range behaviour of a point pattern is governed by
something other than the mechanism that sets its local spacing.

\subsection{An unexplained feature}

The low-$k$ exponent of the count spectrum is small but significantly
negative, $\alpha = -0.14 \pm 0.04$: the plateau turns slightly upward at the
smallest accessible wavevectors. We have no explanation, and we have now
excluded the three candidates we could think of.

It is not a finite-size artefact: Sec.~\ref{sec:finitesize} shows the upturn
surviving a doubling of the box, over which the bias of the estimators
themselves demonstrably disappears. It is not a thresholding artefact: the
exponent moves by less than $0.05$ as the phase cut is swept across the
central $40\%$ of the composition gap. And it is not a large-scale modulation
of domain size---the explanation we regarded as most likely before measuring
it. Were the field to contain regions of systematically larger domains, those
regions would hold fewer of them, and a long-wavelength count fluctuation
would follow; but the mark covariance of Sec.~\ref{sec:attribution} shows no
mass correlation whatever beyond two domain spacings, out to a separation of
$232$ cells, while the excess in $S(k)$ appears at wavelengths several times
longer than that.

What remains is a small excess of long-wavelength count fluctuation in a
pattern whose local order is strong, whose solute is rigid, and whose domain
sizes are spatially uncorrelated on those scales. We report it as an open
question rather than assign it to the nearest available correlate.

\subsection{Limitations}

The analysis is two-dimensional, isothermal, at a single composition, and uses
a single free-energy form. The domain count per configuration is
$\mathcal{O}(10^3)$ and the ensemble is four realisations, so the accessible
wavevector range spans about $1.7$ decades; we therefore report exponents with
their uncertainties over that window and do not assign a hyperuniformity
class. No experimental point pattern is analysed here, and the natural test of
the prediction---that a real precipitate population should show a
count spectrum well above its mass-weighted spectrum---remains to be done on
electron-microscope maps.

% =====================================================================
\section{Conclusion}
% =====================================================================

In a phase-separating system with a locally conserved order parameter, the
solute-rich domains are locally more regular than a hard-core gas; the
concentration field is hyperuniform; and the point pattern of domain centroids
is not, on both the structure-factor and number-variance criteria, in a
measurement whose estimators are validated in the same box and whose verdict
is unchanged when the box is doubled. Weighting the same points by their
solute mass recovers the field's behaviour. A variance budget built on a
parameter-free identity shows the solute-mass fluctuation to be a factor of
eight below what independent domains of the observed size distribution would
produce, with the shortfall supplied by a count--size anticorrelation that
deepens monotonically with scale.

What is conserved is not what is counted. Before classifying a conserved
system as hyperuniform or otherwise from the positions of its objects, one
should ask whether the objects carry the conserved quantity in equal shares,
and measure the weighted spectrum when they do not.

\section*{Data availability}

The analysis code, the derived per-field statistics, and the figure scripts
are openly available at \url{https://github.com/Ruqing1963/conserved-domain-hyperuniformity}. In addition to the summary statistics we release
\texttt{domain\_marks\_stage2000.csv}, the complete per-domain table for the
$1024^2$ ensemble analysed here: $10\,942$ domains across four realisations,
each with its realisation index, centroid coordinates in grid units, area in
cells, and solute mass. Every quantity in Sec.~\ref{sec:attribution} and in
Table~\ref{tab:budget} can be recomputed from that file alone, without
repeating the simulations.

% =====================================================================
% BIBLIOGRAPHY -- checked 17 Aug 2026 against published records, except the
% two entries carrying an inline comment (Tomita 1991, Ezoe 2025), which were
% taken from secondary reference lists.
% =====================================================================

\begin{thebibliography}{99}

\bibitem{TorquatoStillinger2003}
S.~Torquato and F.~H. Stillinger,
Local density fluctuations, hyperuniformity, and order metrics,
Phys. Rev. E \textbf{68}, 041113 (2003).

\bibitem{Torquato2018}
S.~Torquato,
Hyperuniform states of matter,
Phys. Rep. \textbf{745}, 1--95 (2018).

\bibitem{Jiao2014}
Y.~Jiao, T.~Lau, H.~Hatzikirou, M.~Meyer-Hermann, J.~C. Corbo, and S.~Torquato,
Avian photoreceptor patterns represent a disordered hyperuniform solution to a
multiscale packing problem,
Phys. Rev. E \textbf{89}, 022721 (2014).

\bibitem{TorquatoWeighted2026}
S.~Torquato, J.~Kim, M.~A. Klatt, R.~Car, and P.~J. Steinhardt,
Hyperuniformity of weighted particle systems,
Phys. Rev. X \textbf{16}, 011042 (2026).

\bibitem{Ezoe2025}
A.~Ezoe, M.~Katori, and T.~Shirai,
Weighted point configurations with hyperuniformity: an ecological example and
models,
J. Phys. Soc. Jpn. \textbf{94}, 064002 (2025).

\bibitem{Chen2026}
R.~Chen,
Pockmark fields are locally ordered but not hyperuniform: a structure-factor
analysis of Mediterranean seabed fluid-escape depressions,
Zenodo (2026), doi:10.5281/zenodo.20854182.

\bibitem{CahnHilliard1958}
J.~W. Cahn and J.~E. Hilliard,
Free energy of a nonuniform system. I. Interfacial free energy,
J. Chem. Phys. \textbf{28}, 258--267 (1958).

\bibitem{Yeung1988}
C.~Yeung,
Scaling and the small-wave-vector limit of the form factor in phase-ordering
dynamics,
Phys. Rev. Lett. \textbf{61}, 1135--1138 (1988).

\bibitem{Furukawa1989}
H.~Furukawa,
Numerical study of multitime scaling in a solid system undergoing phase
separation,
Phys. Rev. B \textbf{40}, 2341--2347 (1989).

\bibitem{Tomita1991}
H.~Tomita,
Preservation of isotropy at the mesoscopic stage of phase separation processes,
Prog. Theor. Phys. \textbf{85}, 47--56 (1991).

\bibitem{DeLuca2024}
F.~De~Luca, X.~Ma, C.~Nardini, and M.~E. Cates,
Hyperuniformity in phase ordering: the roles of activity, noise, and
non-constant mobility,
J. Phys.: Condens. Matter \textbf{36}, 405101 (2024).

\bibitem{Wilken2023}
S.~Wilken, A.~Chaderjian, and O.~A. Saleh,
Spatial organization of phase-separated DNA droplets,
Phys. Rev. X \textbf{13}, 031014 (2023).

\bibitem{ClarkEvans1954}
P.~J. Clark and F.~C. Evans,
Distance to nearest neighbor as a measure of spatial relationships in
populations,
Ecology \textbf{35}, 445--453 (1954).

\end{thebibliography}

\end{document}
